\chapter{Cold Sores (Herpes Labialis)}

\section*{Definition}
Recurrent, painful, vesicular eruptions on the lips and perioral skin caused by herpes simplex virus (typically HSV-1), characterized by prodromal tingling followed by clustered blisters that crust and heal within 7-14 days.

\section*{When to See a Doctor}
See your doctor if:
\begin{itemize}
  \item Frequent recurrent episodes (more than 6 per year), lesions spreading to eyes, or lesions lasting more than 2 weeks
\end{itemize}

\section*{How It Develops}
Primary HSV-1 infection typically occurs in childhood (gingivostomatitis). The virus then establishes lifelong latency in the trigeminal ganglion. Reactivation triggers include UV radiation, fever, stress, immunosuppression, trauma, and hormonal changes. Reactivated virus travels along sensory nerve axons to the skin, causing epithelial cell lysis, vesicle formation, and local inflammation. Healing involves cell-mediated immune clearance.

\section*{Causes}
\begin{itemize}
  \item HSV-1 reactivation (most common, triggered by UV light, stress, illness)
  \item Fever or viral illness (hence ``cold sores'')
  \item Ultraviolet radiation (sun exposure, tanning beds)
  \item Physical or emotional stress
  \item Immunosuppression (HIV, chemotherapy, transplant)
  \item Hormonal changes (menstruation)
  \item Facial trauma or dental procedures
  \item Fatigue and sleep deprivation
  \item Cold weather or wind exposure
  \item Close contact with active lesions (autoinoculation or transmission)
\end{itemize}

\section*{Related Symptoms}
\begin{itemize}
  \item Fever
  \item Swollen lymph nodes
  \item Sore throat
  \item Mouth ulcers
  \item Facial pain
\end{itemize}

\section*{Medical Evaluation}
\begin{itemize}
  \item Your doctor will diagnose cold sores based on the characteristic appearance of vesicular lesions on an erythematous base at the vermilion border of the lips.
  \item Viral culture or PCR testing of vesicular fluid is performed when the diagnosis is uncertain or when HSV-2 involvement is suspected in atypical presentations.
  \item Blood antibody testing usually does not confirm that a current lip lesion is a cold sore or distinguish a first episode from a recurrence; it is reserved for selected situations.
  \item Referral to an infectious disease specialist or immunologist is arranged for severe, frequent, or treatment-resistant recurrences, particularly if immunosuppression is suspected.
\end{itemize}

\section*{Treatment Options}
\begin{itemize}
  \item Topical antiviral creams (acyclovir 5\%, penciclovir 1\%) applied 5 times daily reduce healing time by 1–2 days when started during the prodrome.
  \item Oral antiviral therapy (acyclovir 400 mg, valacyclovir 2 g, or famciclovir 1.5 g) is used for severe episodes or as a single-day course for episodic treatment.
  \item Daily suppressive antiviral therapy may be considered for frequent, severe, or disruptive recurrences after discussing benefits, risks, and local prescribing guidance.
  \item Analgesic mouthwashes or topical anaesthetics (lidocaine 2\%) provide symptomatic relief for painful lesions during eating and speaking.
\end{itemize}

\section*{Self-Care and Home Management}
\begin{itemize}
  \item If using a topical antiviral where available, start it at the first sign of tingling and follow the product instructions; oral antivirals may work better for some people.
  \item Keep the affected area clean and dry, and avoid touching the sores to prevent spreading the virus to other parts of your body or to others.
  \item Use a lip balm with SPF 30+ sunscreen daily to prevent UV-triggered recurrences, and avoid known triggers such as stress and fatigue.
  \item Apply a cold compress to the sores for 10 minutes several times daily to reduce pain and inflammation during the blister phase.
  \item Avoid sharing utensils, cups, towels, lip products, or razors with others while lesions are active.
\end{itemize}

\section*{References}
\begin{thebibliography}{99}
\bibitem{cold_sores:ref1} Whitley RJ, Kimberlin DW, Prober CG. ``Herpes simplex virus infections.'' In: Nelson Textbook of Pediatrics. 22nd ed. Elsevier; 2024.
\bibitem{cold_sores:ref2} Gilbert SC. ``Management and prevention of recurrent herpes labialis.'' Am Fam Physician. 2022;106(4):408-414.
\bibitem{cold_sores:ref3} Chi CC, Wang SH, Delamere FM, et al. ``Interventions for prevention of herpes simplex labialis.'' Cochrane Database Syst Rev. 2015;(8):CD010095.
\bibitem{cold_sores:ref4} Lebrun-Vignes B, Bouvet E, Alcaraz I, et al. ``Management of herpes labialis.'' Ann Dermatol Venereol. 2022;149(1):15-22.
\bibitem{cold_sores:ref5} Sauerbrei A. ``Herpes labialis: an update.'' Hautarzt. 2018;69(6):480-487.
\bibitem{cold_sores:ref6} Wald A. ``Treatment and prevention of herpes simplex virus type 1 infection in adults.'' UpToDate. Wolters Kluwer; 2025.
\end{thebibliography}
